% EJEMPLO REMOVIBLE: sustituya este bloque por contenido validado de la propuesta.
\subsection*{Ejemplo de patrón editorial}
El siguiente bloque demuestra el orden requerido para migrar evidencia visual y referencias. La Figura~\ref{fig:flujo-ejemplo} resume un flujo conceptual entre operación, datos y decisiones; se presenta sólo como ejemplo de composición y no constituye una definición de la solución.

\begin{figure}[htbp]
  \centering
  \begin{tikzpicture}[
  node distance=12mm,
  bloque/.style={rounded corners=2mm,minimum width=34mm,minimum height=14mm,align=center,draw=SynLine,fill=SynSurface,text=SynNavy,font=\bfseries},
  flujo/.style={-{Stealth[length=2.5mm]},line width=1.2pt,draw=SynCyan}
]
  \node[bloque] (operacion) {Operaciones};
  \node[bloque,right=of operacion] (datos) {Datos gobernados};
  \node[bloque,right=of datos] (decision) {Decisiones};
  \draw[flujo] (operacion) -- node[above,font=\fontsize{9}{11}\selectfont,text=SynMuted]{captura} (datos);
  \draw[flujo] (datos) -- node[above,font=\fontsize{9}{11}\selectfont,text=SynMuted]{contexto} (decision);
  \draw[SynTeal,line width=1pt,dashed,-{Stealth[length=2.3mm]}]
    (decision.south) to[bend left=32] node[below,font=\fontsize{9}{11}\selectfont,text=SynMuted]{retroalimentación} (operacion.south);
\end{tikzpicture}

  \caption{Flujo conceptual de información para apoyar decisiones}
  \label{fig:flujo-ejemplo}
  \FuenteFigura{Elaboración propia para demostrar el uso de la plantilla.}
\end{figure}

El diagrama evidencia cómo una figura debe apoyar un argumento previamente introducido. Después de la figura, el texto debe interpretar relaciones, consecuencias y decisiones; no basta con repetir sus etiquetas.

La Tabla~\ref{tab:ejemplo-trazabilidad} muestra una estructura breve de cinco columnas adecuada para resumir trazabilidad. Los detalles extensos deben mantenerse en anexos o formularios independientes.

\begin{table}[htbp]
  \centering
  \caption{Ejemplo compacto de trazabilidad}
  \label{tab:ejemplo-trazabilidad}
  \begin{tabularx}{\textwidth}{@{}YYYYY@{}}
    \toprule
    \EncabezadoTabla{Necesidad} & \EncabezadoTabla{Diseño} & \EncabezadoTabla{Entrega} & \EncabezadoTabla{Prueba} & \EncabezadoTabla{Operación} \\
    \midrule
    Información oportuna & Integración gobernada & Flujo habilitado & Evidencia registrada & Indicador monitoreado \\
    \bottomrule
  \end{tabularx}
  \FuenteFigura{Elaboración propia para demostrar el estilo tabular.}
\end{table}

La tabla permite seguir una necesidad hasta su control operacional. En el documento real, cada fila debe usar identificadores y evidencias consistentes con los demás subdocumentos.

\begin{decisionarquitectura}
Una decisión debe registrar contexto, alternativas evaluadas, fundamento, consecuencias y responsable de aprobación. \textcite{iso27001} puede citarse cuando una decisión se sustente efectivamente en dicho estándar.
\end{decisionarquitectura}

Los antecedentes corporativos detallados deben entregarse mediante la \ReferenciaAnexo{Formulario T-6}{SYNAPTIX-Formulario-T-6.pdf}, sin incrustar el formulario en este PDF.
